\chapter{还有哪些问题没有答案}

\begin{kaichang}
如果你去过海边，不妨做一件小事：装一瓶海水带回家，放在桌上。它看起来清清亮亮，和一瓶普通的矿泉水没什么两样。

现在告诉你关于这瓶水的三个事实。第一，一毫升海水——大约二十滴——里通常住着上百万个细菌，还有上千万个病毒颗粒，比一座大城市的人口还密集。第二，这些居民中的绝大多数，科学家至今叫不出名字：我们知道它们存在，能在 DNA 测序结果里读到它们的“签名”，却从未在实验室里把它们培养出来、看清楚过。第三，也是最大的一个：这瓶水本身就是一桩悬案的案发现场。约四十亿年前，正是在这样的水里（或者类似的水边），完成了从“普通化学物质”到“会自我复制的系统”的那次跨越——而那次跨越具体是怎样发生的，到今天没有人知道答案。

这本书用了八十六章，讲原子怎样搭成分子、分子怎样组成细胞、信息怎样在 DNA 里保存和改写。读到这里，你可能会以为科学已经把这条链上的每个环节都摸清了。最后一章的任务恰恰相反：带你去看一看这条链上仍然黑着的几盏灯。
\end{kaichang}

\section{“不知道”也分好几种}

先说清楚一件事：科学家的“不知道”，和考试时的“不知道”，不是同一种东西。

考试里的不知道，常常是“书上有答案，我没记住”。而科学前沿的不知道，至少分三种。第一种是\textbf{还没去查}：深海某条海沟里住着什么物种，没人采过样，纯粹是还没人去看。第二种是\textbf{正在查、已有候选答案}：有几个互相竞争的假说，每个都能解释一部分证据，大家正在设计实验区分它们——这一章的大部分问题属于这一种。第三种更深：\textbf{我们还不确定该用什么概念去问}。比如“意识是怎么回事”，目前连“怎样测量一个系统有没有意识”都没有公认的办法，问题本身还在被打磨。

这三种“不知道”，你在生活里其实已经见过。医生说“这个症状的原因尚不明确”，可能是还没查到（第一种），也可能是几种可能都在排除中（第二种）。天气预报说“三天后有 60\% 的概率下雨”，不是预报员偷懒，而是大气是一个对初始条件极其敏感的系统：测量差一点点，几天后的结果就完全不同。这种“原则上只能给概率”的不知道，本章后面还会回来找你。

判断一个“不知道”是不是好的科学问题，有三个标准：边界清楚（我们知道哪些部分已经搞懂了）、有候选解释（不是一片空白，而是几个可以区分的猜想）、有下一步（存在能做的实验或观测）。带着这把尺子，我们来看五个至今还开着的大问题。

\section{生命究竟怎样起源}

把镜头推回约四十亿年前的地球。那时的海洋里没有鱼，没有藻，可能连氧气都几乎没有，只有水、溶解的盐、简单的气体，以及火山和闪电送来的能量。今天这瓶海水里的每一个细菌，都是那场变化的后人。

我们已经知道了什么？链条的两端相当清楚。起点一端：简单的无机和有机小分子，在合适的条件下可以自发地组装出生命的“零件”——氨基酸、核苷酸、脂肪酸，这一点被实验和陨石证据反复支持（证据故事马上讲）。终点一端：所有现存生命共享同一套遗传密码、同一批核心分子机器，说明它们来自共同的祖先（第 45、56 章讲过这些证据）。核糖体——第 69 章那台翻译蛋白质的机器——其核心由 RNA 构成，像一块活化石，暗示历史上可能真的存在过一个“RNA 世界”：RNA 分子既当信息载体，又当催化剂，身兼两职。

缺的是中间：从“一堆能反应的小分子”到“一个会复制、会出错、会被自然选择挑拣的系统”，中间究竟经过了几站？发生在深海热泉喷口的矿物微孔里，还是火山脚下时干时湿的浅水池里，还是冰层之中？脂质小泡（第 47 章）先出现，还是 RNA 先出现？这些问题每一个都有认真的候选答案，也都有专门检验它们的实验，但没有一个能宣布胜利。诚实的回答是：\textbf{生命起源的大方向可信，具体路线未知}。

两个主要候选场景，各有一手好牌，也各有软肋。深海热泉派的理由是：喷口旁的矿物微孔像无数间微型反应室，孔壁上的铁硫矿物能催化反应，热水与冷海水之间的温差还能提供持续的能量流——今天最深处的微生物，确实有不少靠这类化学能过活。反对者则指出：海水太多太咸，核苷酸这类零件难以浓缩到能互相碰撞、连成链的浓度。浅池派恰好相反：时干时湿的小水洼能把分子蒸浓、晒干、再冲开，干—湿循环在实验室里真的促成了长链分子的生成；可批评者问：四十亿年前陆地有没有足够多稳定的水洼？紫外线那么强，刚合成的链会不会很快被破坏？你看，每一派都能做实验、都能被反驳、也都还没被驳倒——这就是“正在查、已有候选答案”的那种“不知道”的真实长相。

\begin{zhengju}
1953 年，芝加哥大学一位名叫米勒的研究生，在导师尤里的支持下搭了一套玻璃装置：一个烧瓶装水代表海洋，加热产生“海蒸气”；另一个大烧瓶装着甲烷（\chem{CH_4}）、氨（\chem{NH_3}）和氢气（\chem{H_2}）的混合气体，代表早期大气；中间用电火花模拟闪电。循环一周后，水里出现了好几种氨基酸——蛋白质的零件，从一锅“无机汤”里自己长了出来。这个结果轰动一时：它第一次证明，“从化学到生命的第一步”是可以搬进实验室检验的。

但故事的下半部分同样重要。后来的研究重新评估了早期大气的成分：对火山气体的分析表明，原始大气里甲烷和氢气可能没有那么丰富，二氧化碳（\chem{CO_2}）和氮气可能更多。在这样的气体里重复米勒式实验，产量确实下降了——原实验的一个前提被修正了。可科学家没有停在这里：他们换配方、换能源（紫外线、冲击波、热泉矿物表面），发现在许多种合理的早期环境下，氨基酸等零件照样能够生成。与此同时，陨石分析补上了另一条路：1969 年落在澳大利亚的默奇森陨石里，检测出七十多种氨基酸，其中不少是地球生命不使用、实验室污染也解释不了的——也就是说，有机零件在太空里就能造出来，甚至可能被陨石“快递”到早期地球。

这个故事的看点不在“米勒答对了”，而在方法：假说被提出，前提被质疑，实验被改进，新的证据渠道被打开。七十年过去，“零件从哪里来”基本有了答案，而“零件怎样组装成会复制的系统”仍然没有——这正是科学正常的呼吸节奏。
\end{zhengju}

\begin{shiyan}
\textbf{巴斯德式对照实验：汤会自己“生出”生命吗？}

材料：两份相同的米汤（或煮过少量蔬菜的清水，放凉），两个干净、带盖的透明杯子，保鲜膜，记号笔。

步骤：① 两份汤分别装入杯子，都装到七分满。② A 杯不盖盖，只用保鲜膜松松罩住、戳几个小孔，让空气能进出；B 杯盖紧盖子密封。③ 并排放在室内明亮但不暴晒的地方，每天同一时间隔着杯壁观察一次，记录浑浊度、表面有没有膜、颜色变化，连续观察五到七天。注意：只许看，\textbf{不要开盖，不要凑近闻}。

观察点：哪一杯先变浑？两杯的变化方式一样吗？如果敞口杯先变浑，它支持“浑浊来自外界落入的东西”，还是“汤自己产生了生命”？再想想巴斯德的鹅颈瓶（第 45 章讲过这个故事）：鹅颈让空气进得去，却挡住了微生物——他改动的到底是哪个变量？

安全提示：变浑的汤里繁殖了未知微生物。实验全程\textbf{不开盖、不品尝}；结束后整杯密封丢进垃圾桶，杯子用开水烫洗，然后洗手。不要用奶制品（腐败气味更重），也不要往汤里加任何“促生长”的东西。
\end{shiyan}

\section{折叠、细胞和意识之间隔着几层楼}

第二个没有答案的问题，藏在尺度里。这本书的因果链是一条不断“上楼”的链：原子组成分子，分子组成细胞，细胞组成组织，组织组成大脑。可是每上一层楼，我们都会发现：自己并不真的会“从楼下的规则，算出楼上的行为”。楼梯断了好几级。

先看蛋白质这一级。第 48 章讲过，蛋白质是一条氨基酸链，必须折叠成精确的三维形状才能工作。这里有一个著名的数字难题，叫\emph{列文塔尔佯谬}。

\begin{qiaoliang}
取一条含 100 个氨基酸的短蛋白质。假设每个氨基酸在链上只有 3 种可能的摆法（真实远不止），那么整条链可能的折叠形状大约有 $3^{100}$ 种。估一下这个数：$\lg 3^{100} = 100 \times 0.477 \approx 47.7$，也就是约 $5\times 10^{47}$ 种。假如蛋白质靠瞎试来寻找正确形状，每试一种只要 $10^{-13}$ 秒（分子动一下的量级），试遍所有形状需要约 $5\times 10^{34}$ 秒。而宇宙的年龄只有约 $4\times 10^{17}$ 秒。换句话说，靠瞎试，一亿亿亿亿个宇宙寿命都不够。可是你的细胞里，真实蛋白质折叠只要千分之一秒到几秒——\textbf{差了三十多个数量级}。这说明折叠根本不是随机搜索，它一定有被能量地形“引导”的路径。这条引导路径长什么样、能不能从序列直接算出来，就是蛋白质折叠问题。近年来，人工智能程序（如 AlphaFold）已经能根据序列相当准确地\emph{预测}折叠好的静态形状，给结构生物学和药物设计帮了大忙；但“链在细胞里、在分子伴侣协助下、一边合成一边折叠的\emph{过程}究竟是怎样的”，仍远远没有解决。预测了终点，不等于画出了路线。
\end{qiaoliang}

折叠之上，还有更陡的楼梯。一个细胞里，上千种蛋白质、脂质和小分子在同时反应（第 55 章的代谢网络），这些化学反应怎样涌现出“细胞会自己决定分裂还是等待”这样的行为？第 57 章只开了个头。再往上：你的大脑约有 860 亿个神经元，每个都只是服从离子通道和神经递质的化学规律（第 58 章），可这么多细胞放在一起，产生了“你正在读这句话、并且理解它”这件事。意识——主观体验本身——怎样从物质组织里长出来？我们不仅没有答案，还在争论该用什么量去度量它。

这里难到什么程度？连“判据”都缺。麻醉医生能让你可靠地失去意识、再安全地醒来，靠的是经验加药理，但“麻醉药究竟关掉了什么”至今没有统一解释；医生判断昏迷病人残存多少意识，只能靠几套间接的脑部测量，而这些测量背后的理论仍在互相竞争。一边是分子层面每个零件都认得，一边是“体验”这个层面连尺子都还在造——两层之间隔着什么，就是我们说的断梯。

这不是说中间的楼层永远不可知，而是说：\textbf{从下层规则预言上层整体行为，是今天科学最吃力的活儿}。这本书反复追问的第二个贯穿问题——“组成部分怎样连接和排列”——在这里变成了最难的形式。

\section{基因组里的暗地图}

第三个问题回到 DNA。你的基因组大约有 30 亿个字母，可是直接编码蛋白质的部分只占 1.5\% 左右——大约两万个基因的“正文”。剩下的 98\% 以上是什么？

第 70、71 章已经告诉你，那里绝不只是垃圾：有开关基因的启动子和增强子，有制造各种调控 RNA 的区段，有染色体的结构件，也有大量来历不明的重复序列——其中一部分是古代病毒留下的“化石”。但“其中多少有功能、有什么功能、有多重要”，仍然是一幅画了一半的地图。

这幅地图正在当众修改，而这正是它特别值得一看的原因。2012 年，一个叫 ENCODE 的大型国际合作项目宣布：人类基因组 80\% 以上的区域都能检测到某种“生化活性”，比如被转录成 RNA。许多媒体立刻报道：“垃圾 DNA 不存在了！”但许多进化生物学家立刻反驳：细胞转录一段序列，可能只是机器偶尔空转；\textbf{“能测到活性”不等于“有功能”}。判断真正的功能，要问第 76 章教你的那个问题：如果这段序列突变掉，生物会受影响吗？自然选择在乎它吗？按这个更严格的标准估计，有功能的序列可能只占基因组的一到两成。这场争论到今天没有完全收场——而它恰好示范了本章的主题：一个响亮的宣布，被同行用更严格的标准检验、压低、修正。科学的“知道”，就是这样一寸一寸挣来的。

\section{比天气预报还难：从基因型到表型}

第四个问题，你可能已经猜到了——第 73 章整章都在为它铺垫：知道了一个人的全部基因序列，能不能预测他会成为什么样的人？

答案是：能预测一点点，而且常常只是概率。身高这样的性状，有成百上千个基因位点各贡献几毫米，再加上营养和童年经历，今天的统计模型已经能解释人群身高差异的一大部分。但轮到心脏病、糖尿病、抑郁这类复杂疾病，基因给出的只是“风险比平均值高或低一些”，环境、生活方式和纯粹的发育随机性都掺在里面（第 85 章讲基因检测时反复强调过：风险不等于诊断）。至于性格、天赋和人生选择，基因的预测力弱到基本可以忽略。

为什么这么难？因为这条因果链要穿过太多层带反馈的楼：基因调控蛋白质，蛋白质塑造细胞，细胞建成大脑，大脑还要回应环境；而环境的一部分——你的家庭、学校、朋友——又会反过来影响基因的打开和关闭（第 74 章）。最好的证据来自同卵双胞胎：他们的 DNA 序列几乎完全相同，常常在同一屋檐下长大，可他们仍然会一个过敏一个不过敏，一个内向一个外向，甚至患上不同的疾病。连“序列相同、环境相近”这种最有利的条件下，预测都会落空，可见发育过程里掺进了多少偶然——哪次细胞分裂出了小错、哪个神经元恰好连上了哪个，都会留下终身印记。这和天气是同一类麻烦：变量太多，彼此纠缠，对初始条件敏感。而且人比大气更麻烦：大气分子不读报纸，你会。任何针对你的预测，一旦被你知道，就可能改变你的行为，让预测失效。从基因型、环境与发育的共同作用算出表型，是生物学眼下最想攻下的预测难题——它既是科学问题，也正在变成每个人都要面对的伦理问题。

\section{把问题倒过来：怎样设计未来}

前四个问题都是“读懂自然”。第五个问题是把这本书倒过来用：\textbf{会预测，才会设计}。我们能不能设计出可持续的材料、药物、能源和农业系统？

要完成的任务，这本书都埋过线索。能源：植物用叶绿素和阳光固定二氧化碳（第 54 章），效率其实不高，但胜在零排放、能自我复制——人工光合作用的梦想，就是用催化剂直接拿阳光、水和 \chem{CO_2} 造燃料，难点在于找到便宜又耐久的催化剂，第 30 章的学问在这里变得值钱。农业：今天养活近一半人类的氮肥，来自一百多年前的哈伯法，它每年用掉全世界百分之几的能源；而豆科植物根瘤里的细菌，在常温常压下干着同样的事。学会模仿甚至改进生物固氮，是农业化学最大的悬赏之一。材料：塑料污染（第 43 章）的出路之一，是设计“用完能被微生物拆开”的高分子，这需要把聚合物化学和代谢网络一起读懂。医药：折叠预测加上基因序列（第 83 章），让“为一个突变定制药物”第一次显得现实。

这些问题的共同结构是：目标在上层（清洁的能源、够吃的粮食、治得好的病），手段在下层（催化剂、序列、网络）。把上下层接通，正是这本书教你的那套问法。它也最诚实地提醒我们：技术带来的不只是能力，还有第 84 章讨论过的外溢与双重用途风险——会设计的人，必须同时会问“该不该”。

\begin{wuqu}
“既然科学自己都承认还没搞明白，那说明各种说法都差不多有道理——顺势疗法、水有记忆，说不定将来会被证明是对的。”这个推理错在把两种“不知道”混为一谈。一种是\textbf{还没有答案}：比如生命起源的具体路线，几种假说都活着，证据还不足以判决。另一种是\textbf{已经有答案，而且答案是否定}：顺势疗法“稀释到分子都不剩还能治病”的说法，被高质量试验反复检验，效果不比安慰剂好；“水记得溶过什么”也与物理化学的全部证据冲突（第 23 章）。前沿的未知，不是给已被否掉的主张开后门的许可证。判断标准始终是那一条：你的主张能不能被检验？检验结果如何？“还不知道 X 怎样发生”，绝不等于“关于 X 的任何故事都一样可信”。
\end{wuqu}

\section{那瓶海水，再问一遍}

回到开场那瓶海水。现在用贯穿全书的五个问题，逐一向它开火：

\begin{itemize}[itemsep=2pt]
\item 它由什么组成？——水、盐、上百万个细菌和病毒，其中大部分没有名字。\textbf{开着}。
\item 组成部分怎样连接和排列？——这些微生物怎样组成群落、谁吃谁、谁给谁发化学信号，第 86 章刚开了头，大部分未知。\textbf{开着}。
\item 什么能量和概率规则决定它能否变化？——这一条已经相当清楚：热力学和动力学规则，与烧杯里那套完全一样。\textbf{基本关上}。
\item 变化有多快、怎样受控制？——海洋吸收 \chem{CO_2} 的速度、海水酸化的速率、碳循环的反馈，有模型，但精度有限。\textbf{半开着}。
\item 在生命中它怎样被保存、复制、调节和选择？——这滴水里的每个生命都携带 DNA，复制、调控、进化样样俱全；可这套系统\emph{最初}怎样启动，正是本章第一问。\textbf{源头开着}。
\end{itemize}

你看，五个问题不是一张考完就扔的问卷，而是一套能带上路的家什。任何新东西摆到面前——一则“突破性疗法”的新闻、一种新材料、一条“科学家发现”的推送——先把这五问依次递上去，你就已经超过了大多数读者。

\section{这本书合上了，问题还开着}

八十六章之前，这本书从厨房里的一粒盐出发，沿着原子、分子、反应、细胞、基因、进化，一路走到生态系统和基因编辑。这一章故意停在五盏还黑着的灯下。这不是泄气，而是这本书最想让你带走的东西：科学的信用，从来不来自“我们已经知道一切”，而来自它知道自己哪里不知道，并且知道怎样把“不知道”变成“知道”——提出可以被检验的假说，让证据来判决，判错了就改。周期表预言了镓，米勒的烧瓶被修正了前提，ENCODE 的 80\% 正在被重新审判：全书每一个证据故事，都是同一套呼吸。

往后的路标，这本书也给你留下了。符号和工具忘了，去翻附录 A、B、C；名词卡住了，附录 D 是词典；下次再看到“研究表明”四个字，先用附录 E 的清单过一遍；练一练做不下去，附录 F 有提示和更远的读物。而比附录更重要的路标，是你已经带走的那五个问题。它们不属于任何一本书，属于任何一个愿意追问的人。

很多年后，你也许会在某个实验室、某块田里、某间诊室里，亲手碰到这一章的某个问题；也许不会。但无论你在哪里，有一件事始终是真的，也是这本书想留给你的最后一句话：

你、树叶、药片和遥远恒星由同一套元素组成；生命的奇妙，不在于它逃离了化学规律，而在于普通分子竟能组织起来，保存过去的信息，并创造仍会改变的未来。

\section*{练一练}

\begin{sikaoti}
\item 选一则最近的科学新闻，用五个贯穿问题逐一向它提问，并标出：哪些问题新闻给了答案，哪些没有，哪些新闻假装有答案其实没有。把结果整理成一张五栏表。
\item 为“生命起源于深海热泉”和“生命起源于时干时湿的浅池”两个假说，各设计一条能区分它们的观测或实验（可以在地球上做，也可以往火星或木星的卫星上想）。说明：什么样的结果会支持其中一个、打击另一个？
\item 列文塔尔佯谬说折叠不是瞎试。请画一张“能量地形”示意图：横轴是折叠的进展，纵轴是能量，用一条起伏下落的通道表示被引导的折叠路径，再在旁边画一条“随机搜索”的对照线。图旁注明：横轴其实代表天文数字种构象，这张图只是帮助思考的表示法。
\item 你的基因组里编码蛋白质的部分只有约 1.5\%。一位朋友说：“剩下 98.5\% 是垃圾，删了也没事。”另一位说：“剩下 98.5\% 全是宝，一个都不能少。”用 ENCODE 之争教你的区分（“有活性”与“有功能”），写三句话，分别回应这两个人。
\item “天气预报都不准，所以科学家的话都不可信。”请用本章对三种“不知道”的区分反驳这句话，要求指出：天气预报的误差属于哪一种“不知道”，它和“已被证据否掉”的主张有什么本质区别。
\item 全书结束前的最后一题：从你身边任选一样物品（一杯水、一片叶子、一粒药……），用五个贯穿问题写一份“已知与未知清单”——每个问题下，各写一条“科学已经知道的”和一条“科学还不知道的”。写完你会发现：你已经有能力把任何一样东西，接回那条从原子通往进化的链上。
\end{sikaoti}
